\section{The \texttt{SAXSCurve} McXtrace Component}
A component mimicking the scattering from a given I(q)-curve by using
linear interpolation between the given points.

\subsection*{Identification}
\begin{itemize}
  \item \textbf{Author:} Martin Cramer Pedersen (mcpe@nbi.dk)
  \item \textbf{Origin:} KU-Science
  \item \textbf{Date:} May 2, 2012
\end{itemize}

\subsection*{Description}
A box-shaped component simulating the scattering from any given I(q)-curve. The component uses linear interpolation to generate the points necessary to compute the scattering of any given photon.

\subsection*{Input parameters}
Parameters in \textbf{boldface} are required; the others are optional.

\begin{longtable}{p{0.22\textwidth}p{0.12\textwidth}p{0.46\textwidth}p{0.14\textwidth}}
\toprule
\textbf{Name} & \textbf{Unit} & \textbf{Description} & \textbf{Default} \\
\midrule
\endhead
DeltaRho & cm/\AA{}$^{3}$ & Excess scattering length density of the particles. & 1.0e-14 \\
Volume & \AA{}$^{3}$ & Volume of the particles. & 10000.0 \\
Concentration & mM & Concentration of sample. & 0.01 \\
AbsorptionCrosssection & 1/m & Absorption cross section of the sample. & 0.0 \\
\textbf{xwidth} & m & Dimension of component in the x-direction. &  \\
\textbf{yheight} & m & Dimension of component in the y-direction. &  \\
\textbf{zdepth} & m & Dimension of component in the z-direction. &  \\
\textbf{SampleToDetectorDistance} & m & Distance from sample to detector (for focusing the scattered x-rays). &  \\
\textbf{DetectorRadius} & m & Radius of the detector (for focusing the scattered x-rays). &  \\
FileWithCurve & str & Datafile with the given I(q). & "Curve.dat" \\
\bottomrule
\end{longtable}

\subsection*{Links}
\begin{itemize}
  \item Component source code found in file \texttt{SAXSCurve.comp}.
\end{itemize}
\IfFileExists{contrib/SAXSCurve_static.tex}{\input{contrib/SAXSCurve_static.tex}}{}